\chapter*{Presentación}
\markboth{Presentación}{Presentación}
\addcontentsline{toc}{chapter}{Presentación}

Este temario propone un repaso riguroso de análisis real de una variable. Parte de la completitud de los números reales, desarrolla sucesiones y series, y continúa con la topología de la recta, continuidad, derivación, convergencia de funciones e integración de Riemann. El objetivo es comprender las definiciones y demostrar los teoremas centrales, además de reconocer cuándo se aplican.

Stephen Abbott, \emph{Understanding Analysis}, segunda edición, es la guía principal de lectura: acompaña de manera progresiva la construcción de los conceptos y sus pruebas. Rudin funciona como una segunda lectura más compacta y exigente. Spivak y Apostol aportan ejercicios y una perspectiva rigurosa del cálculo diferencial e integral \cite{abbott2015,rudin1976,spivak2008,apostol1974}.

La lista de capítulos conserva la división propuesta para estas notas. Los temas de sucesiones de funciones, teorema fundamental del cálculo y Taylor conectan los resultados de convergencia con las herramientas clásicas del cálculo. La integración se plantea inicialmente en el sentido de Riemann, con énfasis en criterios de integrabilidad, propiedades e intercambio justificado de límites.
